\section{Cluster-label evaluation protocol}
\label{app:baseline-labels}

We evaluate labeling over a 28-node PACS hierarchy slice (four fields, each with two categories and two subcategories). 
All methods are given the same cluster members and asked to assign scientific field labels. Compared methods include \doctolora, \texttt{ICAE}, \texttt{KeyLLM}, \texttt{vec2text}, \texttt{T2L} (three embedding spaces), and a baseline (\texttt{in-context}) using only the base Qwen3-4B LLM prompted with abstracts from the 20 members nearest each centroid. \texttt{ActPatch} (hidden-state baseline) is evaluated separately in Appendix~\ref{app:actpatch}.

\doctolora receives its instruction with the cluster documents internalized in the adapter and no document text in the prompt:
\begin{quote}\itshape
In 2 to 3 words, name the scientific field that all of these documents belong to.
Reply with only the field name.
\end{quote}
\texttt{ICAE} receives the same prompt verbatim and decodes against its averaged memory slots.
\texttt{T2L} receives the same prompt and decodes using its training prompt format.
\texttt{KeyLLM} and \texttt{vec2text} cannot take this instruction.
\texttt{KeyLLM} is a keyword extractor returning a ranked list of keywords.
\texttt{vec2text} is an iterative inverter trained to reconstruct text from an embedding and has no prompt interface.
We therefore score every method on its native output without rewriting any label through a second language model.
Passing outputs through an auxiliary namer would test the namer.
The comparison asks whether the representation itself supports the abstraction.

\begin{table}[h]
    \centering
    % The "Operates on" column carries a long cell (the three T2L spaces), which
    % pushed the tabular ~31pt past \textwidth at \small. Two steps down, with
    % tighter column padding, brings it inside the margin without a \resizebox.
    \scriptsize
    \setlength{\tabcolsep}{4pt}
    \caption{Cluster-labeling protocols comparing embedding-arithmetic (top) and text-based (bottom) methods without downstream language model post-processing: ``Operates on'' denotes the representation averaged and labeled, ``\#docs'' indicates consumed member documents per node, and ``Native output'' defines directly evaluated raw outputs.}
    \label{tab:label-protocols}
    \begin{tabular}{@{}l l l l@{}}
    \toprule
    Method & Operates on & \#docs & Native output (scored as-is) \\
    \midrule
    \doctolora & embedding (mean adapter) & up to $2{,}000$ & field name (2--3 words) \\
    \texttt{ICAE} & embedding (mean memory slots) & $8$ & field name (2--3 words) \\
    \texttt{T2L} & embedding (\texttt{GTE} input, task encoder, or LoRA factors) & up to $2{,}000$ & field name (2--3 words) \\
    \texttt{vec2text} & embedding (mean GTR vector, inverted) & $20$ & reconstruction (no prompt) \\
    \midrule
    \texttt{KeyLLM} & text (medoid document) & $1$ & keyword list \\
    \texttt{in-context} & text (centroid-nearest abstracts) & $20$ & field name (2--3 words) \\
    \bottomrule
    \end{tabular}
\end{table}

This design categorizes the methods based on output length. Compared to the official PACS labels of $4.0$ words, both \doctolora and \texttt{ICAE} respond to prompts with $2.4$ words each. Therefore, their scores are directly comparable. \texttt{KeyLLM} averages $8.8$ words, while \texttt{vec2text} averages $15.5$ words. The inability to adhere to a concise naming format is an inherent characteristic of the methods, and we do not apply length corrections. \texttt{T2L} outputs average $2.0$ words. Across all methods, \doctolora, \texttt{ICAE}, \texttt{T2L}, and \texttt{vec2text} operate exclusively on embeddings, whereas \texttt{KeyLLM} and the \texttt{in-context} read the member text. Among the embedding methods, \doctolora differs in that it reads the abstraction level from the length of the mean adapter vector (Appendix~\ref{app:length-dial}).

\textbf{Fuzzy overlap} compares native outputs directly against official PACS labels using a deterministic token-set ratio in \texttt{rapidfuzz}.
The score ranges from $0$ to $1$, involves no language model, and reproduces deterministically from the released strings.
Because fuzzy overlap measures surface tokens, verbose outputs receive lower scores against four-word reference labels, and the metric cannot credit valid paraphrases.

The \textbf{pairwise win rate} evalutes each method using five language-model judges (Reproducibility Statement, p.~\pageref{sec:repro}). For each cluster, judges compare the official PACS label and two candidate method labels, picking the closer one in meaning or declaring a tie; all pairs are shown in both label orders to remove position bias, i.e., a win counts only if both orders agree, otherwise scored as a tie ($0.5$). 
Eight methods enter the round robin: three \texttt{T2L} variants, four other methods, and the \texttt{in-context}. Reported win rates are averaged over all nodes and shown in \figref{cluster-labels}b as row-vs-column cells. \texttt{T2L} is shown only for its latent space that best performs on this task. 
Judges rate \doctolora labels as tied with official PACS names on $32\%$ of judge-node decisions.
By comparison, this rate is $25\%$ for \texttt{ICAE}, $22\%$ for the \texttt{in-context}, and $0\%$ for both \texttt{KeyLLM} and \texttt{vec2text}.

% Generated by the label_eval pipeline (workflow/rules/baseline_trees.smk): fuzzy overlap and
% word counts from label_eval_summary.md, pairwise from label_eval_metric4.json (`snakemake label_eval`).
% Every method is scored on its native output (label_eval_prep.py METHOD_FILES); no naming LLM.
\begin{table}[h]
\centering
\small
\caption{
    Cluster-label quality over $28$ PACS nodes. Each method's native output is scored against the true node label. \textbf{Fuzzy overlap:} deterministic token-set ratio with the reference. \textbf{Pairwise:} win rate from a five-judge round robin over eight methods (five methods, three \texttt{T2L} spaces, and \texttt{in-context}); ties score $0.5$ and require both label orders to agree. $\pm$ is the bootstrap s.d.\ over $1{,}000$ node resamples. \textbf{Words:} mean label length (\textit{PACS:} $4.0$). The ground-truth control has no win rate. Best (excluding control) in bold. Full pairwise matrix: \figref{cluster-labels}b.
}
\label{tab:label-eval}
\setlength{\tabcolsep}{5pt}
\begin{tabular}{@{}l c c c@{}}
\toprule
Method & Words & Fuzzy overlap & Pairwise \\
\midrule
\textit{Ground truth (control)} & $4.0$ & $1.000$ & --- \\
\midrule
\doctolora & $2.4$ & $\mathbf{.642 \pm .055}$ & $\mathbf{.847 \pm .019}$ \\
\texttt{ICAE} & $2.4$ & $.573 \pm .058$ & $.791 \pm .024$ \\
\texttt{KeyLLM} & $8.8$ & $.311 \pm .019$ & $.671 \pm .026$ \\
\texttt{vec2text} & $15.5$ & $.302 \pm .026$ & $.462 \pm .016$ \\
\texttt{T2L} (best of $3$ spaces) & $2.0$ & $.277 \pm .010$ & $.149 \pm .004$ \\
\midrule
\texttt{in-context} & $2.4$ & $.499 \pm .051$ & $.794 \pm .020$ \\
\bottomrule
\end{tabular}
\end{table}

\doctolora achieved the highest scores on both metrics, with a fuzzy overlap of $.642 \pm .055$ and a pairwise win rate of $.847 \pm .019$ (\tabref{label-eval}).
All five judges ranked \doctolora first, with each judge’s win rate ranging from $.696$ to $.786$.
\doctolora’s direct baseline is \texttt{ICAE}, which processes identical instructions and generates outputs of the same average length.
\doctolora outperforms \texttt{ICAE} in pairwise win rates ($.847$ vs. $.791$).
Since their marginal intervals overlap in terms of fuzzy overlap ($.642 \pm .055$ vs. $.573 \pm .058$), we examine the paired node-level differences.
The mean paired difference is $+.069$, the bootstrap standard deviation is $.034$, and the $95\%$ confidence interval is $[.016, .145]$.
Out of a total of $28$ nodes, \doctolora achieved a higher overlap at $12$ nodes, \texttt{ICAE} achieved it at $1$ node, and $15$ nodes were tied.
In direct head-to-head comparisons, \doctolora and \texttt{ICAE} tied in $103$ out of $140$ decisions; \doctolora won approximately two-thirds of the remaining decisions, resulting in a head-to-head win rate of $.596$.
Compared to \texttt{in-context} (pairwise $.794$, fuzzy $.499$), \doctolora recorded a head-to-head win rate of $.629$, with $78$ ties.
Therefore, decoding an averaged adapter recovers the field at least as well as reading the member abstracts.

\doctolora is the top performer on both metrics, consistently ranked first by all judges. 
Compared to its closest baseline, \texttt{ICAE}, which operates under the same conditions, \doctolora leads in win rate, and their other scores are similar enough to warrant direct comparison. Analysis at the node level shows \doctolora typically outperforms or ties with \texttt{ICAE}. 
In head-to-head matchups, \doctolora has a clear edge, and it also surpasses the \texttt{in-context} approach. 
Overall, decoding an averaged adapter is at least as effective as reading member abstracts.

\texttt{T2L} performs worst on both metrics, outputting just a few repeated labels (e.g., ``Social psychology'') across all $28$ nodes, with a fuzzy overlap floor of $.277$. Its three embedding spaces are nearly indistinguishable and lose almost every pairwise comparison, including against \texttt{vec2text}. The method's mapping does not yield decodable representations. The synthesized adapters expanded from the frozen \texttt{GTE} encoder cannot generate meaningful labels. We note that cluster-labeling is not the intended use case and goal for \texttt{T2L}.

Native outputs for every PACS node appear in \tabref{hierarchy-labels} in Appendix~\ref{app:tables}.
\doctolora outputs match official designations, producing ``Quantum chromodynamics'' for code 12.38 and ``Fluid dynamics'' for code 47.
\texttt{KeyLLM} reflects single-document specificity, returning ``square-lattice Heisenberg antiferromagnet, S=1/2, nearest-neighbor exchange'' for magnetic ordering (75.10).
\texttt{ICAE} identifies ``Condensed matter physics'' for code 71 and ``Nonlinear optics'' for code 42.65, but occasionally shifts hierarchy levels, returning ``Materials science'' for bulk band structure (71.20) and ``Quantum many-body systems'' for general physics (0).
\texttt{vec2text} reconstructions are garbled, yielding ``in-distance interactions between electron-bonding spectral properties'' for code 3 and ``inferromagnetic spin-diabetic scattering interactions'' for code 75.

\subsection{Incoherent-cluster control}
\label{app:incoherent}

We test whether decoded cluster labels require a coherent member set (\tabref{incoherent}).
For each of $30$ real PACS nodes, we construct two control clusters of identical size by sampling papers uniformly across PACS chapters outside the target node.
We decode one control at its native centroid length of $11.15$ and a second control scaled to the real-node mean of $11.49$.

A five-judge panel sees each label alongside twelve sampled member titles and is not told which method or cluster type a row comes from.
The panel accepts the real-node label on $73.3\%$ of clusters and rejects $3.3\%$ (one of $30$).
The remainder are judged partial.
Judges accept zero control labels, rejecting $100\%$ of the native control and $96.7\%$ of the norm-matched control.
Control decodings collapse almost uniformly to ``Quantum field theory''.
The norm-matched control earns the same verdict as the native control.


\begin{table}[h]
\centering
\footnotesize
\caption{Scores from the blinded five-judge panel on the incoherent-cluster control experiment. \emph{PACS node (real)} denotes actual clusters, while the control rows consist of clusters of the same size sampled across different PACS chapters and decoded either at their native centroid norm or rescaled to the real nodes' norm. \emph{belongs=yes} and \emph{rejected} indicate the proportion of clusters whose label the panel accepted as valid for the members or rejected entirely. The rest were judged partial.}
\label{tab:incoherent}
% auto-generated by workflow/scripts/groupc/incoh_score.py (issue #101) -- do not edit
\begin{tabular}{lrrr}
\toprule
cluster set & $n$ & belongs=yes & rejected\\
\midrule
PACS node (real) & 30 & 73.3\% & 3.3\% \\
cross-chapter control & 60 & 0\% & 100\% \\
cross-chapter, norm-matched & 60 & 0\% & 96.7\% \\
\bottomrule
\end{tabular}

\end{table}

\subsection{Adapter length sets label granularity}
\label{app:length-dial}

Vector length governs the level of abstraction in decoded \doctolora embeddings.
\secref{cluster-labeling} shows that scaling an embedding by $\alpha$ shifts the decoded label from specific topics to general fields (for instance, moving from citric acid cycle to cellular respiration) until near the origin the label reverts to the prior distribution of the prompt.
We evaluate this magnitude sweep on \texttt{ICAE} and \texttt{vec2text} across five seed documents covering a physics concept, a biochemistry process, a fruit, a historical person, and a patent.
\figref{cluster-labels}d shows the \doctolora sweep on four of these seeds.
The person seed appears only in the tables of Appendix~\ref{app:tables}.

Neither baseline exhibits upward semantic abstraction under vector scaling.
\texttt{ICAE} remains semantically unchanged across lengths, returning the specific concept label at every value of $\alpha$ with variations limited to capitalization and surface wording.
\texttt{vec2text} maintains surface-level specificity while the reconstruction remains coherent, then decomposes into unrelated text fragments at low $\alpha$ without naming broader fields (\tabref{radial-vec2text}).
Among the three embedding methods on this sweep, the specific-to-general ladder is specific to \doctolora.
